\chapter{胶子算符、极化投影与重正化基底}
\label{sec:b05-gluon-operators}

本章把“胶子分布”拆成三个不能互相替代的层次：连续场论中的非局域算符、格点上的有限距离矩阵元，以及把格点量映射到光锥分布的方案链。胶子场属于伴随表示，场强又带两个反对称 Lorentz 指标，因此同一个物理分布在写进格点代码之前必须先固定色表示、指标缩并、Wilson 线方向与重正化基底。

\section{双胶子算符与 Lorentz 不变振幅}

\paragraph{定义.}
沿空间方向 \(z^\mu=z\hat z^\mu\) 连接两个场强的规范协变算符写成
\begin{equation}
  O_g^{\mu\nu\rho\sigma}(x;z)
  =F_a^{\mu\nu}(x+z\hat z)\,
   U_z^{ab}(x+z,x)\,
   F_b^{\rho\sigma}(x),
  \label{eq:b05-gluon-operator}
\end{equation}
其中 \(U_z\) 是基础表示链接构造的平行移动；在矩阵元中通过
\begin{equation}
  U_z^{ab}(x+z,x)
  =2\,\operatorname{Tr}\!\left[T^aU_z(x+z,x)T^bU_z(x,x+z)\right]
  \label{eq:b05-adjoint-link}
\end{equation}
解释其伴随指标。每个场强可选取 \(6\) 个独立 \(( \mu,\nu)\) 分量，因此双线性一般需要 \(6\times6=36\) 个 Lorentz 结构来审计。

\paragraph{公式：六个不变振幅.}
对自旋平均的外部态，定义矩阵元
\begin{equation}
  M_{\mu\alpha;\lambda\beta}(z,p)
  =\langle p|G_{\mu\alpha}(z)W[z,0]G_{\lambda\beta}(0)|p\rangle .
  \label{eq:b05-bare-matrix}
\end{equation}
Lorentz 协变性把它分解为
\begin{align}
  M_{\mu\alpha;\lambda\beta}
  &=(g_{\mu\lambda}p_\alpha p_\beta
    -g_{\mu\beta}p_\alpha p_\lambda
    -g_{\alpha\lambda}p_\mu p_\beta
    +g_{\alpha\beta}p_\mu p_\lambda)M_{pp}\nonumber\\
  &\quad+(g_{\mu\lambda}z_\alpha z_\beta
    -g_{\mu\beta}z_\alpha z_\lambda
    -g_{\alpha\lambda}z_\mu z_\beta
    +g_{\alpha\beta}z_\mu z_\lambda)M_{zz}\nonumber\\
  &\quad+(g_{\mu\lambda}z_\alpha p_\beta
    -g_{\mu\beta}z_\alpha p_\lambda
    -g_{\alpha\lambda}z_\mu p_\beta
    +g_{\alpha\beta}z_\mu p_\lambda)M_{zp}\nonumber\\
  &\quad+(g_{\mu\lambda}p_\alpha z_\beta
    -g_{\mu\beta}p_\alpha z_\lambda
    -g_{\alpha\lambda}p_\mu z_\beta
    +g_{\alpha\beta}p_\mu z_\lambda)M_{pz}\nonumber\\
  &\quad+(p_\mu z_\alpha-p_\alpha z_\mu)
         (p_\lambda z_\beta-p_\beta z_\lambda)M_{ppzz}\nonumber\\
  &\quad+(g_{\mu\lambda}g_{\alpha\beta}
    -g_{\mu\beta}g_{\alpha\lambda})M_{gg}.
  \label{eq:b05-six-amplitudes}
\end{align}
六个振幅依赖于 \(z^2\) 和 Ioffe 时间 \(\nu=-(p\cdot z)\)。

时间反演与场交换使 \(M_{pp},M_{zz},M_{gg},M_{ppzz}\) 以及
\(M_{pz}-M_{zp}\) 为 \(\nu\) 的偶函数，而 \(M_{pz}+M_{zp}\) 为奇函数。
这个奇偶结构是后面用 \(z\leftrightarrow -z\) 分离极化信号的场论根据。

\paragraph{公式：非极化投影.}
在光锥极限 \(z=z^-\hat n_-\) 下，最低扭度胶子分布来自 \(M_{pp}\)：
\begin{equation}
  -M_{pp}(\nu,z^2=0)
  =\frac12\int_{-1}^{1}dx\,e^{-ix\nu}x f_g(x),
  \qquad f_g(x)=g(x).
  \label{eq:b05-mpp-to-pdf}
\end{equation}
注意被积函数含 \(x\)：这里得到的是动量密度 \(xg(x)\) 的 Fourier 变换，而不是 \(g(x)\) 本身。实验和全局拟合若使用不同的 \(xg(x)\) 约定，必须统一后才能比较。

\section{非极化、螺旋度与 Ioffe-时间投影}

格点上不能直接令 \(z^2=0\)，只能选取一组在同圈具有相同反常量纲的等时分量。常用的非极化组合为
\begin{equation}
  \mathcal O_{\mathrm{unpol}}(z)
  =M_{tx;tx}(z)+M_{ty;ty}(z)-2M_{xy;xy}(z).
  \label{eq:b05-unpol-combination}
\end{equation}
在部分有效场论讨论中可写成
\begin{equation}
  M_{ti;it}+M_{ij;ji}=2p_0^2M_{pp}+M_{\mathrm{tr}}+\cdots ,
  \label{eq:b05-pp-combination}
\end{equation}
其中 \(M_{\mathrm{tr}}\) 概括 \(M_{gg}\)、迹缩并和高扭度接触项。\(M_{ij;ji}\) 并不能在所有约定下被当成 \(M_{pp}\) 的无污染替代；式 \eqref{eq:b05-pp-combination} 的余项必须在算子基底和匹配方案中保留。

\paragraph{极化投影.}
引入对偶场强
\begin{equation}
  \widetilde G^{\mu\nu}=\frac12\epsilon^{\mu\nu\rho\sigma}G_{\rho\sigma},
  \qquad \epsilon^{0123}=+1,
  \label{eq:b05-dual-field}
\end{equation}
并定义极化矩阵元
\begin{equation}
  \widetilde m_{\mu\alpha;\lambda\beta}(z,p,s)
  =\langle p,s|G_{\mu\alpha}(z)W[z,0]\widetilde G_{\lambda\beta}(0)|p,s\rangle .
  \label{eq:b05-polarized-matrix}
\end{equation}
自旋无关的背景由偶部进入；螺旋度相关部分用
\begin{equation}
  \widetilde M_{\mu\alpha;\lambda\beta}(z,p,s)
  =\widetilde m_{\mu\alpha;\lambda\beta}(z,p,s)
  -\widetilde m_{\mu\alpha;\lambda\beta}(-z,p,s)
  \label{eq:b05-zodd}
\end{equation}
提取。格点上应同时计算 \(+z\) 与 \(-z\)，并保持链接路径反向、端点场强顺序和组态样本完全相同。

极化 Ioffe-时间分布由
\begin{equation}
  -i\widetilde I_p(\nu,z^2=0)
  =\widetilde M_{(ps)}^{(+)}(\nu,0)-\nu\widetilde M_{(pp)}(\nu,0)
  =\frac{i}{2}\int_{-1}^{1}dx\,e^{-ix\nu}x\Delta g(x)
  \label{eq:b05-polarized-ioffe}
\end{equation}
给出。线性组合中减去 \(\nu\widetilde M_{(pp)}\) 的作用，是消去由外部动量与自旋结构产生的非螺旋度振幅，而不是对数据作经验性减背景。

\paragraph{约化赝 ITD.}
令同一极化投影的坐标矩阵元为 \(\widetilde M(z,p)\)。定义
\begin{equation}
  \mathfrak M(\nu,z^2)
  =\frac{\widetilde M(z,p)}
         {\widetilde M(0,p)},
  \label{eq:b05-reduced-itd}
\end{equation}
并与微扰核 \(K\) 匹配：
\begin{equation}
  \mathfrak M(\nu,z^2)
  =\int_{0}^{1}du\,K(u\nu,\mu^2z^2)I_p(u\nu,\mu)
  +\mathcal O(z^2\Lambda_{\mathrm{QCD}}^2).
  \label{eq:b05-pseudo-factorization}
\end{equation}
比值可消去只依赖 \(z\) 的乘性 UV 因子，但不能消去硬核、软函数、rapidity 方案或大 \(z\) 幂修正。因此“约化后无显式 \(Z\)”不等于“无需重正化论证”。

\begin{table}[H]
  \centering
  \caption{胶子投影的用途与主要污染。}
  \label{tab:b05-projections}
  \small
  \begin{tabularx}{\textwidth}{@{}p{0.23\textwidth}p{0.25\textwidth}X@{}}
    \toprule
    投影 & 目标量 & 必须报告的余项或几何 \\
    \midrule
    \(M_{tx;tx}+M_{ty;ty}-2M_{xy;xy}\) & 非极化 \(xg(x)\) & \(M_{gg}\)、迹项、Wilson 方向、匹配核 \\
    \(\widetilde M_{0i;0i}+\widetilde M_{ij;ij}\) & 螺旋度 \(\Delta g(x)\) & 对偶张量符号、\(z\)-奇组合、核子自旋投影 \\
    约化矩阵元 \(\mathfrak M\) & 赝 ITD & 归一化点、短距窗口、\(z^2\) 幂修正 \\
    \(F W F\) 的任意单项 & 诊断量 & 一般不是 twist-2 的本征算符，不能单独发布 PDF \\
    \bottomrule
  \end{tabularx}
\end{table}

\section{可乘性重正化、三角混合与构建块}

\paragraph{严格的纯准算符陈述.}
在规范不变的 UV 正规化中，纯准胶子算符的端点和 Wilson 线自能可由有限个因子组织。按四个场强指标中含 \(z\) 方向的个数 \(s\) 分类，可把结构写成
\begin{equation}
  O_{g,R}^{\mu\nu\rho\sigma}(z)
  =e^{-\delta m|z|}Z_{\mathrm{op}}^{-1}(s)
   O_{g,B}^{\mu\nu\rho\sigma}(z),
  \label{eq:b05-multiplicative}
\end{equation}
\begin{equation}
  Z_{\mathrm{op}}(s)
  =Z_{wg}Z_{vg1}^{s/2}Z_{vg2}^{(2-s)/2}.
  \label{eq:b05-zop}
\end{equation}
这里 \(s=0,1,2\)；\(\delta m\) 概括格点正规化下的线性自能，\(Z_{wg}\) 是 Wilson 线端点因子，\(Z_{vg1},Z_{vg2}\) 区分含 \(z\) 指标的分量。由于整体 \(Z\) 依赖 \(s\)，任意改变各分量权重的线性组合未必保持可乘性。

\paragraph{辅助场闭包中的三角混合.}
为了系统识别小基底内的重定义，可引入伴随表示辅助场
\begin{equation}
  \mathcal L_{\mathrm{aux}}
  =\mathcal L_{\mathrm{QCD}}
  +\overline Q_a(x)(in\cdot D-m)_{ab}Q_b(x).
  \label{eq:b05-auxiliary-lagrangian}
\end{equation}
在固定规范与维数正规化中，与场强--辅助场双线性可能重定义的有三组结构：
\begin{equation}
  \begin{pmatrix}J_{1,R}^{\mu\nu}\\J_{2,R}^{\mu\nu}\\J_{3,R}^{\mu\nu}\end{pmatrix}
  =
  \begin{pmatrix}
    Z_{11}&Z_{12}&Z_{13}\\
    0&Z_{22}&Z_{23}\\
    0&0&Z_{33}
  \end{pmatrix}
  \begin{pmatrix}J_1^{\mu\nu}\\J_2^{\mu\nu}\\J_3^{\mu\nu}\end{pmatrix}.
  \label{eq:b05-triangular-mixing}
\end{equation}
其中 \(J_2\) 是场强与辅助场组成的规范不变结构，\(J_3\) 比例于辅助场运动方程。物理矩阵元中 \(J_3\) 的接触项可消失，但在检查格点算子基底时不能把它的存在从代数分类中删掉。

\paragraph{四个构建块.}
积分掉辅助场后，可作为独立重正化单元的结构为
\begin{align}
  \mathcal O_R^1
  &=Z_{11}^2e^{\delta m\Delta z}
    F^{ti}W_{z}F^{ti},\label{eq:b05-block1}\\
  \mathcal O_R^2
  &=Z_{22}^2e^{\delta m\Delta z}
    F^{zi}W_{z}F^{zi},\label{eq:b05-block2}\\
  \mathcal O_R^3
  &=Z_{11}Z_{22}e^{\delta m\Delta z}
    F^{ti}W_{z}F^{zi},\label{eq:b05-block3}\\
  \mathcal O_R^4
  &=Z_{22}^2e^{\delta m\Delta z}
    F^{z\mu}W_{z}F^{z}_{\ \mu},
  \label{eq:b05-block4}
\end{align}
其中横向指标 \(i\) 需按声明的方式缩并，\(\Delta z\) 是两端纵向坐标差。指数因子表达的是自能结构；实际格点拟合仍要用多格距和同几何 Wilson 圈检验。

\begin{table}[H]
  \centering
  \caption{胶子算符重正化中两类“混合”的区别。}
  \label{tab:b05-two-mixings}
  \small
  \begin{tabularx}{\textwidth}{@{}p{0.22\textwidth}p{0.28\textwidth}X@{}}
    \toprule
    层次 & 数学对象 & 物理含义与边界 \\
    \midrule
    UV 算符基底 & \(O_{i,R}=Z_{ij}O_j\) & 同量子数、同量纲的格点算符重定义；决定能否用标量因子或必须用矩阵 \\
    辅助场闭包 & \(J_{1,2,3}\) 三角矩阵 & 用于分类端点和 Wilson 线结构，不表示任意物理投影都相互混合 \\
    味单态匹配 & \(C_{gg},C_{gq},C_{qg},C_{qq}\) & 准分布匹配到光锥 PDF 时的胶子--夸克卷积，不是 UV 自能 \\
    投影误差 & \(M_{\mathrm{tr}}\)、高扭度、接触项 & 由选错 Lorentz 组合或省略迹项产生，不能被重正化因子吸收 \\
    \bottomrule
  \end{tabularx}
\end{table}

一个安全的结论是：在指定纯准算符基底和规范不变正规化下，36 个分量可按 \(s\) 分类处理；在包含辅助场或不同物理投影的更大闭包中，仍可能出现三角或完整矩阵混合。把这两句话合并成“36 个分量在所有基底下都无混合”是错误的。

\section{Clover 场强、Wilson 链接与非连通矩阵元}

格点上用四个定向 plaquette 的 Clover 平均构造场强。取围绕同一点的定向 plaquette \(P_{\mu\nu}\)，定义
\begin{equation}
  Q_{\mu\nu}=P_{\mu\nu}+P_{\nu,-\mu}+P_{-\mu,-\nu}+P_{-\nu,\mu},
  \label{eq:b05-clover-sum}
\end{equation}
并以反厄米部分构造
\begin{equation}
  F_{\mu\nu}^{\mathrm{clov}}
  =-\frac{i}{8a^2}\left(Q_{\mu\nu}-Q_{\mu\nu}^{\dagger}\right).
  \label{eq:b05-clover-field}
\end{equation}
不同快照对四个 plaquette 的书写顺序和反向链接定向不完全一致；代码审计应以“同起点、同闭合方向、反厄米”作为归一化条件，而不能只看某一项的显式正负号。

基础表示的有限距离矩阵元写成
\begin{equation}
  M_{\mu\lambda;\nu\rho}(z)
  =\sum_{\vec x}\operatorname{Tr}\!\left[
    F_{\mu\lambda}(x+z\hat z)U_z(x+z,x)
    F_{\nu\rho}(x)U_z(x,x+z)\right],
  \label{eq:b05-lattice-correlator}
\end{equation}
其中第二条链接是返回路径。非极化和螺旋度流分别为
\begin{align}
  \mathcal O_{\mathrm{unpol}}^{(0)}(z)
  &=\left[-F^{0i}W_zF_{0i}+2F^{12}W_zF_{12}\right]_{\mathrm{E}} ,
  \label{eq:b05-lattice-unpol}\\
  \mathcal O_{\mathrm{hel}}^{(0)}(z)
  &=-\Bigl\{F^{01}W_zF^{23}-F^{02}W_zF^{13}
    +2F^{12}W_zF^{03}-(z\leftrightarrow -z)\Bigr\}_{\mathrm{E}} .
  \label{eq:b05-lattice-hel}
\end{align}
第一式中的相对负号来自时间型与空间型场强的 Wick 转动；第二式中的 \(z\)-奇反对称化是螺旋度投影，不是普通的统计对称化。

\paragraph{非连通三点函数.}
胶子流是纯规范算符，在 Wick 收缩后与核子夸克线拓扑分离。逐组态计算时使用
\begin{equation}
  C_{3}(z,P_z;t_s,t_i)
  =\left\langle
    \left(\mathcal O_g(z,t_i)-\langle\mathcal O_g(z,t_i)\rangle\right)
    \left(C_2(P_z;t_s,t_0)-\langle C_2(P_z;t_s,t_0)\rangle\right)
  \right\rangle ,
  \label{eq:b05-disconnected}
\end{equation}
再以 \(C_3/C_2\) 和源汇间隔拟合提取基态矩阵元。真空减除必须在组态层面完成；先对乘积求平均再减去两个均值会保留交叉项的偏置。

\paragraph{验收条件.}
本节的算符构造只有在以下条件同时满足时才可进入匹配阶段：色表示为伴随或显式转换为基础；场强为反厄米且四叶定向一致；Wilson 线方向、端点、长度在裸量和减除量中相同；投影含迹项处理；\(+z\) 与 \(-z\) 数据使用同一源汇与统计样本；非连通真空减除可复现。

\chapter{TMD、准 TMD 与软结构}
\label{sec:b05-tmd}

TMD 不是“在 PDF 上多写一个 \(k_T\)”。它把纵向动量、横向分离、UV 重正化标度和 rapidity 标度组织成同一个因子化对象。格点计算只能直接得到 Euclidean 矩阵元；软函数、CS 核、硬匹配和极限外推必须逐层闭合。

\section{定义、共线极限与 TMD 分类}

在 \(b_T\) 空间定义带 Dirac 投影的裸关联函数
\begin{equation}
  h_{\Gamma}(z,b_T,L,P_z;1/a)
  =\langle P|\bar\psi(0)\Gamma
    W_{\sqsubset}(0,z;b_T)\psi(z,b_T)|P\rangle ,
  \label{eq:b05-bare-tmd}
\end{equation}
其中 \(W_{\sqsubset}\) 是有限长度 staple，\(L\) 是纵向臂长。把 UV 减除后的矩阵元记为 \(h_R\)，准 TMD 的 Fourier 变换写成
\begin{equation}
  \widetilde f_{\Gamma}(x,b_T,P_z)
  =\mathcal N^{-1}\int\frac{dz}{2\pi}\,
    e^{-ixzP_z}h_R(z,b_T,P_z),
  \label{eq:b05-quasi-tmd-fourier}
\end{equation}
\(\mathcal N\) 固定场强归一化和 Fourier 约定。坐标空间截断、大 \(z\) 尾部和 \(x\) 空间非物理区都属于同一反演问题，不能用同一条光滑曲线同时“拟合掉”。

\paragraph{共线极限.}
在 leading twist 且 rapidity 方案固定时，
\begin{equation}
  f_{g/h}(x,\mu)
  =\int d^2k_T\,f_1^g(x,k_T;\mu,\zeta),
  \label{eq:b05-collinear-limit}
  \end{equation}
并满足总动量求和
\begin{equation}
  \int_0^1dx\,x\left[f_g(x,\mu)+\sum_qf_q(x,\mu)\right]=1.
  \label{eq:b05-momentum-sum}
\end{equation}
这两个关系是方案一致性检验。横向积分截断、单态混合、零模减除或重正化误差都可能让数值不满足式 \eqref{eq:b05-collinear-limit}，不能只把它们归入“统计误差”。

\begin{table}[H]
  \centering
  \caption{八种 leading-twist quark TMD 的分类与格点状态。表中状态是 B05 文档所述先例层级，不代表多格距物理闭合。}
  \label{tab:b05-quark-tmd}
  \footnotesize
  \begin{tabularx}{\textwidth}{@{}p{0.18\textwidth}p{0.11\textwidth}p{0.10\textwidth}X@{}}
    \toprule
    函数 & 核子/夸克极化 & T 性质 & 计算状态与关键风险 \\
    \midrule
    \(f_1\) & U/U & 偶 & 核子单系综先例；连续极限与软化未闭合 \\
    \(g_{1L}\) & L/L & 偶 & 需极化核子态；完整胶子对应尚未完成 \\
    \(g_{1T}\) & T/L & 偶 & worm-gear；自旋投影与非连通噪声 \\
    \(h_1\) & T/T & 偶 & 手征奇数；格点混合与端点行为 \\
    \(h_{1L}^{\perp}\) & L/T & 偶 & worm-gear L；需横向极化投影 \\
    \(h_{1T}^{\perp}\) & T/T & 偶 & pretzelosity；高扭度污染风险 \\
    \(f_{1T}^{\perp}\) & T/U & 奇 & Sivers；依赖 Wilson 线方向，仍属探索 \\
    \(h_1^{\perp}\) & U/T & 奇 & Boer--Mulders；核子与介子已有先例，胶子对应未完成 \\
    \bottomrule
  \end{tabularx}
\end{table}

胶子 TMD 不能把表 \ref{tab:b05-quark-tmd} 中的 Gamma 矩阵替换为场强后直接复用。胶子 Wilson 线在伴随表示，场强投影有无迹、对偶和线性极化结构，T-odd 与过程方向的关系也不同。

\section{Staple 几何、四类发散与 Wilson 圈减除}

一个常用的 Euclidean staple 路径为
\begin{equation}
  \Gamma:\quad 0\longrightarrow L\hat z
  \longrightarrow L\hat z+b_T
  \longrightarrow z\hat z+b_T .
  \label{eq:b05-staple-path}
\end{equation}
路径方向和长度是算符定义的一部分。非极化 leading-twist 分布在某些方案中可对方向差异重吸收，但 T-odd、自旋相关和软函数矩阵元不能把 \(+L\) 与 \(-L\) 的裸量混为一谈。

\begin{table}[H]
  \centering
  \caption{staple 关联函数的主要发散与错误处置。}
  \label{tab:b05-staple-divergences}
  \small
  \begin{tabularx}{\textwidth}{@{}p{0.20\textwidth}p{0.25\textwidth}X@{}}
    \toprule
    结构 & 典型形式 & 处置与常见错误 \\
    \midrule
    线性自能 & \(e^{-\delta m(2L+z)/a}\) & 用同长度 Wilson 圈平方根或质量重正化；不同表示不能互借系数 \\
    Cusp & 由 cusp 反常维数控制的对数 & 转角和表示必须同几何；不能用普适常数替代 \\
    Pinch pole & \(e^{-V(b_T)L}\) & 由两条平行长 Wilson 线交换产生；需要 \(L\) 平台或同几何圈减除 \\
    端点对数 & \(\ln a\) 型复合算符发散 & 由短距比率或匹配吸收；普通 RI/MOM 单除不够 \\
    Rapidity & 依赖光锥方向和软胶子 & 属于软函数/rapidity 方案，不是普通 UV 发散 \\
    \bottomrule
  \end{tabularx}
\end{table}

设矩形 Wilson 圈为 \(Z_E(2L+z,b_T;1/a)\)，其渐近结构可概括为
\begin{equation}
  \ln Z_E
  =-2\delta m L-2V(b_T)L
   -\Gamma_{\mathrm{cusp}}\ln(\mu L)
   +C_E(b_T;\mu)+\cdots .
  \label{eq:b05-wilson-loop-asymptotic}
\end{equation}
若开 Wilson 线的线性自能与 pinch 项按总长度增长，则
\begin{equation}
  h^{\mathrm{sub}}(z,b_T,P_z;1/a)
  =\lim_{L\to\infty}
   \frac{h^{\mathrm{bare}}(z,b_T,L,P_z;1/a)}
        {\sqrt{Z_E(2L+z,b_T;1/a)}}
  \label{eq:b05-wilson-loop-subtraction}
\end{equation}
可以抵消主项。式 \eqref{eq:b05-wilson-loop-asymptotic} 是拟合与诊断模板，不是对所有 \(L\) 的严格等式；只有同表示、同转角、同端点和同横向分离的数据才能进入同一平台拟合。

\paragraph{短距比率.}
选取短距参考点 \((z_0,b_{T,0})\)，把剩余端点对数匹配到 \(\overline{\mathrm{MS}}\)：
\begin{equation}
  h^{\mathrm{SDR}}(z,b_T,P_z;\mu)
  =\frac{h^{\mathrm{sub}}(z,b_T,P_z;1/a)}
         {h^{\mathrm{sub}}(z_0,b_{T,0},0;1/a)}
   \frac{h^{\overline{\mathrm{MS}}}(z_0,b_{T,0};\mu)}
        {h^{\overline{\mathrm{MS}}}(z,b_T;\mu)}.
  \label{eq:b05-sdr}
\end{equation}
参考点必须同时满足 \(a\ll z_0,b_{T,0}\ll 1/\Lambda_{\mathrm{QCD}}\)。固定参考点在不同格距上对应不同的物理距离会产生离散化偏差；固定物理参考点又会改变微扰精度，因此参考点误差必须通过扫描或匹配转换显式估计。

\paragraph{RI/MOM 的边界.}
RI/MOM 可以确定某些局域反常量纲和混合矩阵，却不自动消除非局域 Wilson 线的线性自能、pinch pole 或 cusp 结构。B05 文档记录的五格距检验表明，单独采用 RI/MOM 后仍有随 \(z\) 或 \(a\) 增强的格距依赖。可审计结论是“RI/MOM 单独不足，需要与同几何减除联合”，而不是“所有 TMD 只能采用一种方案”。

\section{软函数、Collins--Soper 演化与匹配}

表示 \(R\) 的软函数可写成
\begin{equation}
  S_R(b_T;\mu,Y_1,Y_2)
  =S_{I,R}(b_T;\mu)
   \exp\!\left[K_R(b_T;\mu)\mathcal Y(Y_1,Y_2)\right],
  \label{eq:b05-soft-factorization}
\end{equation}
其中 \(S_I\) 是 rapidity 无关的内禀部分，\(K_R\) 是 Collins--Soper 核，\(\mathcal Y\) 依赖正规化方案。对胶子，\(R=\mathrm{adj}\)，领头 cusp 系数与 \(C_A\) 成正比；quark 的 \(C_F\) 数值不能直接替换。

TMD 的 rapidity 演化和 CS 核的 UV 演化分别为
\begin{equation}
  2\zeta\frac{\partial}{\partial\zeta}\ln f^R(x,b_T;\mu,\zeta)
  =K_R(b_T;\mu),
  \qquad
  \mu^2\frac{\partial}{\partial\mu^2}K_R(b_T;\mu)
  =-\Gamma_{\mathrm{cusp}}^R(\alpha_s).
  \label{eq:b05-cs-evolution}
\end{equation}
软函数在因子化式中以平方根进入，是因为一条 TMD Wilson 线与软函数的两条 Wilson 线几何不同。该平方根、符号和常数项都属于方案，必须与匹配核同步。

\paragraph{结构因子化式.}
准 TMD 到光锥 TMD 的关系可组织为
\begin{equation}
  \begin{aligned}
  \frac{\widetilde f(x,b_T;\zeta_z,\mu)}
       {\sqrt{S_I(b_T;\mu)}}
  =H\!\left(\frac{\zeta_z}{\mu^2}\right)
   \exp\!\left[\frac12\ln\!\left(\frac{\zeta_z}{\zeta}\right)K(b_T;\mu)\right]
   f(x,b_T;\mu,\zeta)\\
  +\mathcal O\!\left(\frac{\Lambda_{\mathrm{QCD}}^2}{\zeta_z},
   \frac{M^2}{P_z^2},\frac{1}{b_T^2\zeta_z}\right),
  \end{aligned}
  \label{eq:b05-tmd-factorization}
\end{equation}
其中 \(\zeta_z\) 由准 TMD 的动量标度确定。不同文献可能使用 \(H\)、卷积或归一化不同的写法；在数值合并前必须统一到同一 Fourier 变换和场强归一化。

\paragraph{硬匹配与大对数.}
固定阶硬核含 \(\ln(\zeta_z/\mu^2)\) 的幂次。在 \(x\) 很小或动量很大时，\(\zeta_z\) 与 \(\mu\) 的比值可能远离 1，固定阶展开失效。重整化群重求和把硬核从自然标度 \(\mu_0\sim 2xP_z\) 演化到参考标度 \(\mu\)。当 \(\mu_0\) 进入 Landau 区时，有限 \(P_z\) 无法直接给出更小 \(x\) 的可控微扰匹配，必须报告 \(x_{\min}\) 和模型/外推误差。

\section{格点 TMD 的证据等级}

\begin{table}[H]
  \centering
  \caption{quark TMD 方法先例与直接 gluon TMD 的边界。}
  \label{tab:b05-tmd-evidence}
  \small
  \begin{tabularx}{\textwidth}{@{}p{0.22\textwidth}p{0.25\textwidth}X@{}}
    \toprule
    证据类型 & 已记录内容 & 可支持与不可支持的结论 \\
    \midrule
    Wilson 圈减除 & 五种格距、有限 \(L\) 平台、不同费米子作用量 & 支持 \(L\to\infty\) 自能减除；不支持完整物理 TMD 已闭合 \\
    \(\pi\) 介子软结构 & 大动量形状因子抽取 \(S_I\)，不同动量比抽取 \(K\) & 支持软核可计算及系综交叉检验；gluon 表示仍需独立匹配 \\
    核子 \(f_1\) & 单系综、NNLO+RGR、核子 TMD 分布先例 & 支持算法链；连续极限和物理体积尚未由同一数据闭合 \\
    Boer--Mulders & 核子与 pion 的手征奇数 T-odd 先例 & 支持 T-odd 方法；不能外推为 gluon T-odd 数值 \\
    直接 gluon TMD & B05 内未见中心值、协方差和多格距组合 & 当前结论只能是框架与工作流，不能发布物理分布 \\
    \bottomrule
  \end{tabularx}
\end{table}

把表 \ref{tab:b05-tmd-evidence} 的不同行相乘、相除或平均会产生错误证据：Wilson 圈五格距验证不能替代核子连续极限；quark 软核不能替代伴随表示软核；单系综 TMD 先例不能替代多格距胶子结果。

\section{梯度流方案与流时验收}

连续梯度流定义为
\begin{equation}
  \partial_tB_\mu(t,x)=D_\nu G_{\nu\mu}(t,x),
  \qquad B_\mu(0,x)=A_\mu(x),
  \label{eq:b05-flow-equation}
\end{equation}
格点上以 Wilson plaquette 作用量离散化。定义流半径
\begin{equation}
  r_f=\sqrt{8t},
  \label{eq:b05-flow-radius}
\end{equation}
并把 \(r_f\) 视为平滑尺度，而不是一个额外的普通重正化标度。固定正流时下，流场规范不变局域关联函数的 UV 行为可由标准 QCD 参数吸收；这不表示任何 Euclidean 非局域矩阵元自动等于光锥物理。

\paragraph{小流时展开.}
短流时下，流化算符按规范不变局域算符展开：
\begin{equation}
  O_R(t,x)
  =\sum_i c_i(t)O_{R,i}(x)
   +\mathcal O(t\Lambda_{\mathrm{QCD}}^2).
  \label{eq:b05-sftx}
\end{equation}
对非局域 TMD，Wilson 系数还必须依赖 \(z,b_T,P_z\)；固定外部物理分离后再对 \(t\) 展开的要求为
\begin{equation}
  a\ll r_f\ll\min\left(|z|,b_T,\Lambda_{\mathrm{QCD}}^{-1}\right),
  \qquad
  P_zr_f\ll1.
  \label{eq:b05-flow-window}
\end{equation}
如果 \(r_f\) 已经接近或超过 \(b_T\)，流会抹平待测的横向结构；如果 \(P_zr_f\) 不小，boost 会把流半径放大成硬尺度；如果 \(r_f\) 太小，格距效应和 SFTX 高阶项又不受控。

\paragraph{flow 不替代的对象.}
梯度流可改善 UV 行为和统计噪声，但不替代 rapidity 重正化、软函数、CS 核、Wilson 线方向或味单态匹配。一个经流化后有限、跨组态稳定的矩阵元仍可能因违反式 \eqref{eq:b05-flow-window}、缺少软减除或使用错误表示而只能作为诊断量。

\paragraph{验证协议.}
至少应做到：扫描多个满足窗口的 \(t\)；在每个 \(t\) 上重新完成 Wilson 圈、场强和准 TMD；用同一匹配核转换到共同方案；对 \(t\to0\) 作受控外推；比较不同外推窗口和匹配阶数；把流时误差、格距、体积、统计和匹配误差分开报告。单一流时的有限结果不足以支持物理分布的中心值。

\chapter{质子自旋、胶子极化与连续极限}
\label{sec:b05-spin-continuum}

自旋问题的困难不在“是否还有别的贡献”，而在不同分解、标度和算符定义下“贡献”指什么。连续极限的困难也不在“多拟合几个参数”，而在格距、动量、距离、流半径和匹配标度之间是否存在非空尺度窗口。

\section{自旋预算与两套分解}

\paragraph{历史量级.}
朴素非相对论夸克模型把质子写成三个 \(S\)-波夸克，容易给出 \(\Delta\Sigma\sim1\)。Ellis--Jaffe 利用 SU(3) 味对称和超子衰变给出
\begin{equation}
  \Delta\Sigma=3F-D\approx0.60\pm0.12,
  \label{eq:b05-ellis-jaffe}
\end{equation}
这已经表明“1”不是唯一自然预期。1987 年 EMC 的测量进一步给出很小的 \(\Delta\Sigma\)，从而把胶子和轨道角动量推到自旋预算的中心。历史数值依赖 QCD 标度、SU(3) 破缺和演化，不能把 EMC 的单一数值直接与现代全局分析逐项相减。

\paragraph{Jaffe--Manohar 分解.}
在无限动量或光锥规范语境中，
\begin{equation}
  \frac12
  =\frac12\Delta\Sigma+L_q+\Delta G+L_G,
  \label{eq:b05-jm}
\end{equation}
其中
\begin{equation}
  \Delta G(Q^2)=\int_0^1dx\,\Delta g(x,Q^2).
  \label{eq:b05-delta-g}
\end{equation}
这个分解能显式分离胶子螺旋度和胶子轨道项，但依赖坐标、规范与过程类。把 \eqref{eq:b05-jm} 中每一项当成与方案无关的局域算符矩阵元，会混淆不同定义。

\paragraph{Ji 分解.}
规范不变的角动量分解写成
\begin{equation}
  J=J_q+J_g,
  \qquad
  J_q=\frac12\Delta\Sigma+L_q,
  \label{eq:b05-ji}
\end{equation}
而 \(J_g\) 是胶子总角动量。\(J_g\) 不能在没有额外方案的情况下再被直接分为“胶子自旋”和“胶子轨道”两部分。两种分解中的 \(\frac12\Delta\Sigma\) 可对照，但 \(L_q,L_G,\Delta G,J_g\) 不能逐个等号连接。

\paragraph{轴反常的作用.}
单味单态轴矢流满足
\begin{equation}
  \partial^\mu A_\mu^0
  =N_f\frac{\alpha_s}{2\pi}\operatorname{Tr}F_{\mu\nu}\widetilde F^{\mu\nu}.
  \label{eq:b05-axial-anomaly}
\end{equation}
反常 Ward 恒等式可以在特定极限下把味单态轴荷与拓扑密度矩阵元联系起来，但反常流不能直接等同于胶子自旋密度。把 \(\Delta\Sigma\) 与 \(\Delta G\) 通过简单加法重复分配，是自旋预算中常见的错误来源。

\section{胶子螺旋度与赝 ITD 证据}

光锥螺旋度分布可由
\begin{equation}
  \Delta g(x,Q^2)
  =\frac{i}{2xP^+}\int\frac{d\xi^-}{2\pi}
    e^{-ixP^+\xi^-}
    \langle P,S|F_a^{+\alpha}(\xi^-)W^{ab}\widetilde F^+_{\alpha,b}(0)|P,S\rangle
  \label{eq:b05-helicity-pdf}
\end{equation}
定义。格点上的极化坐标矩阵元需先做 \(z\)-奇组合，再进入 Ioffe-时间或准分布。HadStruc 的赝 ITD 计算是一个单格距、较重 pion 质量的先例，表 \ref{tab:b05-spin-constraints} 所列数值只能作为来源文献快照。

\begin{table}[H]
  \centering
  \caption{B05 文档记录的 \(\Delta G\) 约束快照。数值未在本任务中独立重分析；比较前须统一标度、积分范围和正定性假设。}
  \label{tab:b05-spin-constraints}
  \small
  \begin{tabularx}{\textwidth}{@{}p{0.25\textwidth}p{0.22\textwidth}X@{}}
    \toprule
    方法 & 来源记录的中心值 & 边界与解读 \\
    \midrule
    DSSV14 & \(0.20^{+0.06}_{-0.07}\) & 全局分析快照，依赖数据集与参数化 \\
    NNPDFpol1.1 & \(0.13\pm0.14\) & 大误差，不能据此声明精确符号 \\
    JAM17 & \(0.06\pm0.21\) & 与零相容，显示约束仍弱 \\
    Ji-scheme 格点 & \(0.251\pm0.047\pm0.016\) & 局域算符方案，匹配系数系统误差较大 \\
    HadStruc 赝 ITD & 仅给出比较趋势 & 单格距探索，不提供完整 \(x\) 依赖或连续极限 \\
    JAM24 & 不要求正值 & 放松正定性后，正负号不能由旧共识外推 \\
    \bottomrule
  \end{tabularx}
\end{table}

对极化结果的验证应优先报告：核子自旋投影；对偶张量符号；\(z\)-奇组合；匹配核是否为 \(\Delta g\) 专用；单态胶子--夸克混合；赝 ITD 的归一化点和 \(z^2\) 范围；以及多格距、物理质量和体积是否闭合。

\section{非极化胶子 PDF 的三格距连续极限}

B05 连续极限主文档记录了三套 CLQCD \(2+1\) 味系综，算符为
\begin{equation}
  \mathcal O_g(z)=M_{tx;tx}+M_{ty;ty}-2M_{xy;xy}.
  \label{eq:b05-continuum-operator}
\end{equation}
Hybrid 方案在短距使用零动量比值，在长距使用自重整化因子：
\begin{equation}
  h_R(z,P_z)
  =\frac{h_B(z,P_z)}{h_B(z,0)}\theta(z_s-|z|)
   +\eta_s\frac{h_B(z,P_z)}{Z_R(z,1/a)}\theta(|z|-z_s).
  \label{eq:b05-hybrid}
\end{equation}
连续性系数 \(\eta_s\) 在 \(z_s\) 处固定重整化后的矩阵元。长距因子采用参数化
\begin{equation}
  Z_R=\exp\left[
    \frac{kz}{a}\ln(a\Lambda_{\mathrm{QCD}})+m_0z
    +\frac{5C_A}{3b_0}\ln\frac{\ln(1/a\Lambda_{\mathrm{QCD}})}
    {\ln(\mu/\Lambda_{\mathrm{QCD}})}
    +\frac12\ln\left(1+\frac{d}{\ln(a\Lambda_{\mathrm{QCD}})}\right)^2
  \right].
  \label{eq:b05-zr}
\end{equation}
这些参数强相关，单独看 \(k\) 或 \(m_0\) 的大误差会误导；应报告 \(Z_R(z)\) 本身及其对最终矩阵元的相关误差。

\begin{table}[H]
  \centering
  \caption{B05 连续极限文档记录的 CLQCD 系综。}
  \label{tab:b05-continuum-ensembles}
  \small
  \begin{tabularx}{\textwidth}{@{}lcccccX@{}}
    \toprule
    系综 & \(a/\mathrm{fm}\) & \(L^3\times T\) & \(m_\pi/\mathrm{MeV}\) & \(m_\pi L\) & \(N_{\mathrm{conf}}\) & 主要用途 \\
    \midrule
    C24P29 & 0.105 & \(24^3\times72\) & 292.3(1.0) & 3.746(13) & 879 & 粗格距与较大体积 \\
    E32P29 & 0.0897 & \(32^3\times64\) & 287.3(2.5) & 4.179(36) & 900 & 中间格距 \\
    F32P30 & 0.0775 & \(32^3\times96\) & 300.4(1.2) & 3.780(15) & 777 & 细格距与高动量检验 \\
    \bottomrule
  \end{tabularx}
\end{table}

表中三套系综的 pion 质量约为 300 MeV，并非物理 pion 点。连续外推不能自动替代手征外推；有限体积参数 \(m_\pi L\) 只覆盖约 3.7--4.2。

\paragraph{大距离外推.}
坐标空间矩阵元在 \(\lambda=zP_z\) 大时信噪比下降。B05 文档采用
\begin{equation}
  h_R(\lambda)=l_1\lambda^{-a_1}e^{-\lambda/\lambda_0}
  \label{eq:b05-lambda-tail}
\end{equation}
补全 Fourier 变换所需尾部。幂律项与 PDF 端点行为相关，指数项概括有限关联长度。该形式必须通过改变拟合区间检验；只用一组区间得到的光滑尾部不能作为独立验证。

\paragraph{联合极限.}
最终外推使用
\begin{equation}
  xg(x,P_z,a)
  =xg_0(x)+a^2f(x)+a^2P_z^2h(x)+\frac{d(x)}{P_z^2}.
  \label{eq:b05-joint-extrapolation}
\end{equation}
\(a^2f\) 是普通离散化误差，\(a^2P_z^2h\) 是 boost 与格距的交叉项，\(d/P_z^2\) 概括高扭度与靶质量修正。三项不能互相吸收；若只有一种 \(P_z\) 或一个格距，式 \eqref{eq:b05-joint-extrapolation} 的系数不可辨识。

\paragraph{结果解读.}
B05 文档报告最终 \(xg(x)/\langle x\rangle\) 在误差内与若干全局拟合一致，大 \(x\) 受压，\(x<0.2\) 与 \(x>0.8\) 被视为 LaMET 不可靠区。PDF 第 6 页的图只支持这一层级；任何具体点的中心值仍应回到原始数据表或发表表格，而不能从扫描图像估读。

\section{误差预算与尺度窗口}

\begin{table}[H]
  \centering
  \caption{胶子 PDF 连续极限的主要误差与可验证构造。}
  \label{tab:b05-continuum-errors}
  \small
  \begin{tabularx}{\textwidth}{@{}p{0.24\textwidth}p{0.25\textwidth}X@{}}
    \toprule
    误差来源 & B05 文档中的构造 & 验收判据 \\
    \midrule
    大 \(\lambda\) 尾部 & 改变拟合起点与终点 & 两组外推在联合误差内一致，且 Fourier 振铃受控 \\
    匹配与标度 & \(\mu=2\to4\ \mathrm{GeV}\) & 标度差进入系统误差，而非常数平移 \\
    连续/无穷动量 & 联合结果与最小格距、较高动量结果比较 & 差值不大于声明的外推误差 \\
    混合方案阈值 & \(z_s=0.3\to0.2\ \mathrm{fm}\) & 中心值变化与 \(z_s\) 相关性显式报告 \\
    统计与拟合 & bootstrap/jackknife、相关协方差 & 重采样整条分析链，不只重采样最后一点 \\
    手征与体积 & 多个 \(m_\pi\)、\(L\) & 本轮 CLQCD 数据只近似 300 MeV，不能宣称已闭合 \\
    \bottomrule
  \end{tabularx}
\end{table}

连续外推还应满足一组可操作窗口：
\begin{equation}
  a\ll |z|,\ b_T\ll \frac{L_s}{2},
  \qquad
  \Lambda_{\mathrm{QCD}}\ll P_z,\ \frac{1}{|z|},
  \qquad
  aP_z\ll1,
  \qquad
  m_\pi L\gtrsim4.
  \label{eq:b05-scale-windows}
\end{equation}
对 flow-TMD 再加式 \eqref{eq:b05-flow-window}。如果窗口为空，正确结果是报告“没有可发布区间”，而不是把窗口外的点纳入同一全局拟合。
\section{从中间量到物理结论的等级}

\begin{table}[H]
  \centering
  \caption{B05 结果的分级验收。}
  \label{tab:b05-claim-levels}
  \small
  \begin{tabularx}{\textwidth}{@{}p{0.21\textwidth}p{0.31\textwidth}X@{}}
    \toprule
    等级 & 已闭合内容 & 允许的表述 \\
    \midrule
    L1：裸格点量 & 场强、链接、矩阵元、真空减除、统计误差 & 可报告诊断量与关联函数 \\
    L2：方案量 & UV 减除、短距参考点、Wilson 圈、比值或 flow 匹配 & 可报告在明确方案下的有限矩阵元 \\
    L3：光锥分布 & soft/rapidity、CS、硬匹配、Fourier、大 \(z\) 尾部 & 可报告 TMD 或 PDF 候选分布及其适用范围 \\
    L4：物理闭合 & 多格距、大动量、手征、体积、单态混合与独立交叉检验 & 才可称为完整的第一性原理物理结果 \\
    \bottomrule
  \end{tabularx}
\end{table}

任何缺少一层的结果都不能用更高层的名称包装。例如，没有软减除的“准 TMD 矩阵元”不是 TMD；单流时的“有限关联函数”不是已匹配的 gluon TMD；单格距的“连续极限比较”不是连续极限。
